% Part: sets-functions-relations
% Chapter: sets
% Section: russells-paradox

\documentclass[../../../include/open-logic-section]{subfiles}


\begin{document}

\olfileid{sfr}{set}{rus}
\olsection{রাসেলের কূটাভাস}

সদস্যভিত্তিক সমতার নীতি অনুযায়ী, $\phi(x)$ ধর্মবিশিষ্ট $x$-গুলির
\emph{সেই নির্দিষ্ট} সেটকে $\Setabs{x}{\phi(x)}$ প্রতীকে লেখা যায়।
তবে এই নীতি \emph{প্রকৃতপক্ষে} কেবল নিচের কথাটিই সমর্থন করে:
\emph{যদি} এমন সেট থাকে, যার সদস্যরা ঠিক সেই সব বস্তু যাদের $\phi$
ধর্মটি আছে, \emph{তবে} এমন সেট একটিই। অন্যভাবে বললে, কোনও~$\phi$
স্থির করা হলে, $\Setabs{x}{\phi(x)}$ সেটটি অনন্য, \emph{যদি সেটি থাকে}।

কিন্তু এই শর্তটি গুরুত্বপূর্ণ!{} সব ধর্মের সাহায্যেই
\emph{ধর্মনির্দেশে সেট গঠন} করা যায় না। অর্থাৎ, কিছু ধর্ম সেট
সংজ্ঞায়িত করে \emph{না}। সব ধর্মই তা করতে পারলে আমরা সরাসরি
স্ববিরোধে পৌঁছতাম। এর সবচেয়ে বিখ্যাত উদাহরণ রাসেলের কূটাভাস।

সেট অন্য সেটের !!{element}s হতে পারে। যেমন, কোনও সেট~$A$-এর ঘাত সেট
সেট দিয়েই গঠিত। কাজেই কোনও সেট অন্য একটি সেটের !!a{element} কি না,
সেই প্রশ্ন করা বা অনুসন্ধান করা অর্থপূর্ণ। একটি সেট কি নিজেরই সদস্য হতে পারে?
সেটের ধারণার মধ্যে এমন কিছু আছে বলে মনে হয় না, যা এই সম্ভাবনা বাদ দেয়।
যেমন, \emph{সমস্ত} সেট যদি বস্তুর একটি সমষ্টি গঠন করে, তবে ভাবা যেতে পারে
যে সেগুলিকে একত্র করে একটি সেট বানানো যায়: সমস্ত সেটের সেট। আর সেটি
নিজেও সেট হওয়ায়, সমস্ত সেটের সেটের !!a{element} হবে।

নিজেকে !!a{element} হিসেবে না রাখার ধর্ম, অর্থাৎ \emph{নিজের সদস্য না হওয়া},
বিবেচনা করলে রাসেলের কূটাভাস দেখা দেয়। যদি ধরে নিই, এমন সমস্ত সেটের
একটি সেট আছে, যারা নিজেরা নিজেদের !!a{element} নয়, তাহলে কী হয়?
সেটটি কি আছে:
\[
R = \Setabs{x}{x \notin x}
\]
আমরা প্রমাণ করতে পারি যে, এমন সেট নেই।

\begin{thm}[রাসেলের কূটাভাস]\ollabel{thm:russells-paradox}
	$R = \Setabs{x}{x \notin x}$ এমন কোনও সেট নেই।
\end{thm}

\begin{proof}
$R = \Setabs{x}{x \notin x}$ সেটটি থাকলে,
$R \in R$ হবে যদি এবং কেবল যদি $R \notin R$ হয়; এটি স্ববিরোধ।
\end{proof}

\begin{tagblock}{novice}
\begin{explain}
প্রমাণটি ধীরে ধীরে দেখি। $R$ থাকলে $R \in R$ কি না, এই প্রশ্ন অর্থপূর্ণ।
ধরি, সত্যিই $R \in R$। কিন্তু $R$-কে সেই সমস্ত সেটের সেট হিসেবে সংজ্ঞায়িত
করা হয়েছে, যারা নিজেদের !!{element}s নয়। তাই $R \in R$ হলে,
$R$ নিজেই $R$-এর সংজ্ঞায় ব্যবহৃত ধর্মটি ধারণ করে না। অথচ কেবল এই ধর্মবিশিষ্ট
সেটগুলিই $R$-এ থাকে। সুতরাং $R$, $R$-এর !!a{element} হতে পারে না,
অর্থাৎ $R \notin R$। কিন্তু $R$ একই সঙ্গে $R$-এর !!a{element} হবে
আবার হবে না, তা অসম্ভব। তাই আমরা স্ববিরোধ পেলাম।

$R \in R$ এই অনুমান থেকে স্ববিরোধ আসে বলে, $R \notin R$।
কিন্তু এ থেকেও স্ববিরোধ আসে!{} কারণ, $R \notin R$ হলে $R$ নিজেই
$R$-এর সংজ্ঞায় ব্যবহৃত ধর্মটি ধারণ করে। তাই নিজেদের সদস্য নয় এমন অন্য
সব সেটের মতো, $R$-ও $R$-এর !!a{element} হবে। আবারও একই সঙ্গে
$R$-এর !!a{element} না হওয়া এবং হওয়া সম্ভব নয়।
\end{explain}
\end{tagblock}

\begin{digress}
রাসেলের কূটাভাস এড়িয়ে কীভাবে সেটতত্ত্ব গড়ে তোলা যায়? অর্থাৎ,
$R = \Setabs{x}{x \notin x}$ সেটটি আছে, এই \emph{অসংগত} দাবি করা
কীভাবে এড়ানো যায়? এর জন্য এমন স্বীকার্য দিতে হবে, যা খুব সুনির্দিষ্টভাবে
বলে দেয় কখন সেটের অস্তিত্ব স্বীকার করা যায়, আর কখন যায় না।

এই অধ্যায়ে সেটতত্ত্বের যে রূপরেখা দেওয়া হয়েছে, তাতে তা করা হয়নি।
এটি \emph{সম্পূর্ণ অনানুষ্ঠানিক}। এখানে শুধু বলা হয়েছে যে, সেটগুলি
সদস্যভিত্তিক সমতার নীতি মেনে চলে এবং কিছু সেট দেওয়া থাকলে তাদের
সংযোগ, ছেদ ইত্যাদি গঠন করা যায়। এর চেয়ে অনেক বেশি কঠোরভাবে
সেটতত্ত্ব গড়ে তোলা সম্ভব। \oliflabeldef{cumul:::part}{সেই কঠোর আলোচনা
\olref[cumul][][]{part} অংশের জন্য রেখে দিচ্ছি। আপাতত আমরা অনানুষ্ঠানিকভাবে
এগোব এবং সতর্কভাবে স্ববিরোধ এড়ানোর চেষ্টা করব।}{}
\end{digress}


\end{document}
